\documentclass[11pt,a4paper]{article}

% Paquetes
\usepackage[utf8]{inputenc}
\usepackage[T1]{fontenc}
\usepackage[spanish]{babel}
\usepackage{geometry}
\usepackage{titlesec}
\usepackage{enumitem}
\usepackage{hyperref}
\usepackage{xcolor}
\usepackage{fancyhdr}
\usepackage{lastpage}

% Configuración de página
\geometry{left=2.5cm, right=2.5cm, top=2cm, bottom=2cm}

% Colores
\definecolor{linkblue}{RGB}{0,102,204}

% Configuración de hipervínculos
\hypersetup{
    colorlinks=true,
    linkcolor=linkblue,
    urlcolor=linkblue,
    pdfauthor={Fabián Belmar},
    pdftitle={Curriculum Vitae - Fabián Belmar}
}

% Formato de secciones
\titleformat{\section}{\large\bfseries\uppercase}{}{0em}{}[\titlerule]
\titlespacing{\section}{0pt}{12pt}{6pt}

\titleformat{\subsection}{\normalsize\bfseries}{}{0em}{}
\titlespacing{\subsection}{0pt}{8pt}{4pt}

% Encabezado y pie de página
\pagestyle{fancy}
\fancyhf{}
\renewcommand{\headrulewidth}{0pt}
\fancyfoot[C]{\thepage\ de \pageref{LastPage}}

% Comandos personalizados
\newcommand{\cventry}[2]{\noindent\textbf{#1}\hfill #2\\}
\newcommand{\cventrysub}[1]{\noindent\hspace{0.5cm}#1\par\vspace{2pt}}

\begin{document}

% ============================================
% ENCABEZADO
% ============================================
\begin{center}
    {\LARGE\bfseries FABIÁN BELMAR}\\[8pt]
    \href{https://orcid.org/0000-0003-4239-1874}{ORCID} |
    \href{https://www.scopus.com/authid/detail.uri?authorId=57220090371}{Scopus} |
    \href{https://www.researchgate.net/profile/Fabian-Belmar-3}{ResearchGate} |
    \href{https://scholar.google.es/citations?user=wr6beBsAAAAJ}{Google Scholar}\\[4pt]
    Correo: \href{mailto:fbelmar@cepchile.cl}{fbelmar@cepchile.cl}\\
    Tel/Whatsapp: +56 9 76371608
\end{center}

\vspace{10pt}

% ============================================
% 1. FORMACIÓN ACADÉMICA
% ============================================
\section{Formación Académica}

\cventry{2020 -- 2025}{Doctorado en Procesos e Instituciones Políticas}
\cventrysub{Universidad Adolfo Ibáñez, Chile. Aprobado con distinción máxima.}
\cventrysub{Tesis: ``Corrupción al margen del Estado. El caso de la FIFA''. Dirigida por Dr. Aldo Mascareño.}

\vspace{6pt}

\cventry{2013 -- 2018}{Licenciado en Ciencia Política y Administración Pública. Administrador Público}
\cventrysub{Universidad de Talca, Chile. Aprobado con distinción máxima.}
\cventrysub{Tesis: ``Vinculación no programática en la gestión municipal. La agenda de los alcaldes, 2015-2016''. Dirigida por Dr. Mauricio Morales.}

% ============================================
% 2. EXPERIENCIA ACADÉMICA
% ============================================
\section{Experiencia en la Academia}

\cventry{2025 --}{Centro de Estudios Públicos (CEP)}
\cventrysub{Investigador asistente del proyecto C22 - Métodos Digitales. Análisis de procesos sociales mediante herramientas de ciencia de datos y técnicas de machine learning.}

\vspace{4pt}

\cventry{2023 -- 2024}{Instituto de Investigaciones Gino Germani, Universidad de Buenos Aires (IIGG)}
\cventrysub{Investigador pasante en el Grupo de Estudios de Culturas Populares Contemporáneas a cargo de la Dra. Verónica Moreira.}

\vspace{4pt}

\cventry{2022 -- 2023}{Proyecto FONDECYT Regular \#1220004}
\cventrysub{Asistente de investigación del proyecto ``Vinculación entre representantes y representados en ambientes formales y virtuales de mediación'' a cargo del Dr. Mauricio Morales.}

\vspace{4pt}

\cventry{2018 -- 2023}{Centro de Análisis Político, Universidad de Talca (CAP-UTALCA)}
\cventrysub{Coordinador de proyectos de investigación y actividades de vinculación con el medio.}

\vspace{4pt}

\cventry{2018 -- 2020}{Proyecto FONDECYT Regular \#1180009}
\cventrysub{Asistente de investigación del proyecto ``Los Partidos Demócrata Cristianos en América Latina'' a cargo del Dr. Mauricio Morales.}

% ============================================
% 3. ARTÍCULOS CIENTÍFICOS
% ============================================
\section{Artículos Científicos}

\subsection{En prensa / En revisión}

\begin{enumerate}[leftmargin=*, label=\arabic*., start=11]
    \item \textbf{Belmar, F.}, \& Mascareño, A. FIFA: Between couplings and hypertrophy. The evolution of the system in the 20th and 21st centuries. \textit{International Journal of Sport Policy and Politics}. [En revisión, 2da ronda]

    \item \textbf{Belmar, F.}, Contreras, G., \& Morales, M. (2025). Who demands clientelism? Differences between individual and organizational demand making. \textit{Political Science Research \& Methods}. [En prensa]

    \item \textbf{Belmar, F.}, \& Mascareño, A. (2025). Corruption: An uneven field of research—Between state and private topics. \textit{Societies}, 15(7), 186. \href{https://doi.org/10.3390/soc15070186}{DOI}

    \item \textbf{Belmar, F.}, Morales, M., \& Villarroel, B. (2025). The Impact of Emotions on Voting Intention in a Constitutional Referendum. \textit{Political Studies Review}. [En prensa]
\end{enumerate}

\subsection{Publicados}

\begin{enumerate}[leftmargin=*, label=\arabic*., start=7]
    \item \textbf{Belmar, F.}, \& Morales, M. (2025). Ethnic solidarity in the non-programmatic distribution of benefits: The case of Mapuche mayors in Chile. \textit{Ethnic and Racial Studies}. \href{https://doi.org/10.1080/01419870.2025.2470896}{DOI}

    \item \textbf{Belmar, F.}, Lara, B., \& Miguieles, J. (2024). Unpacking electoral competition: Campaign spending, vote distribution and turnout in Chilean mayoral elections. \textit{Policy Studies}. \href{https://doi.org/10.1080/01442872.2024.2427280}{DOI}

    \item \textbf{Belmar, F.}, Maldonado, F., \& Morales, M. (2024). In search of the missing link: Election law infractions and candidate sanctions. \textit{Public Integrity}. \href{https://doi.org/10.1080/10999922.2024.2364384}{DOI}

    \item \textbf{Belmar, F.}, Morales, M., \& Villarroel, B. (2023). Writing constitution without parties? The programmatic weakness of party-voter linkages in the Chilean political change. \textit{Politics}, 45(1). \href{https://doi.org/10.1177/02633957231158073}{DOI}

    \item \textbf{Belmar, F.}, Contreras, G., Morales, M., \& Troncoso, C. (2023). Demanding non-programmatic distribution. Evidence from local governments in Chile. \textit{Policy Studies}. \href{https://doi.org/10.1080/01442872.2023.2172718}{DOI}

    \item \textbf{Belmar, F.}, \& Morales, M. (2022). Clientelism, turnout and incumbents' performance in Chilean local government elections. \textit{Social Sciences}, 11(8), 361. \href{https://doi.org/10.3390/socsci11080361}{DOI}

    \item \textbf{Belmar, F.}, \& Morales, M. (2020). Clientelismo en la gestión local. El caso de los alcaldes en Chile, 2015-2016. \textit{Política y Sociedad}, 57(2), 567-591. \href{https://doi.org/10.5209/poso.62315}{DOI}
\end{enumerate}

% ============================================
% 4. CAPÍTULOS DE LIBRO
% ============================================
\section{Capítulos de Libro}

\begin{enumerate}[leftmargin=*, label=\arabic*.]
    \item \textbf{Belmar, F.} (2024). Descentralización y participación ciudadana en Chile: evaluación de la autonomía de los gobiernos locales. En VV.AA. \textit{La Revolución Silente: Descentralización en Chile}.

    \item \textbf{Belmar, F.} (2020). Las tecnologías al servicio de la comunicación y el funcionamiento interno de los parlamentos. En VV.AA. \textit{Parlamentos Conectados}. Universidad de Oviedo, España.

    \item \textbf{Belmar, F.} (2020). La estrategia de comunicación de los representantes: la construcción de la imagen de los políticos. En VV.AA. \textit{Parlamentos Conectados}. Universidad de Oviedo, España.

    \item \textbf{Belmar, F.}, Herrera, M., \& Obregón, M. (2018). El efecto de la elección de concejales sobre la elección de consejeros regionales 2017. En Morales, M., Navarrete, B. \& Vial, C. \textit{Política Subnacional en Chile}. RIL Editores.
\end{enumerate}

% ============================================
% 5. DOCENCIA
% ============================================
\section{Docencia}

\cventry{2025}{Universidad de Talca}
\cventrysub{Evaluación de Políticas, Programas y Proyectos (Cátedra); Ciencia de la Administración Pública (Cátedra).}

\cventry{2025}{Universidad Diego Portales}
\cventrysub{Métodos Cuantitativos Introductorios (Cátedra); Métodos Cuantitativos Avanzados (Cátedra).}

\cventry{2025}{INACAP}
\cventrysub{Introducción a las Políticas Públicas (Cátedra); Sistemas de Compras Públicas (Cátedra).}

\cventry{2023 -- 2024}{Universidad Central}
\cventrysub{Análisis de Datos Cuantitativos (Cátedra); Métodos Cuantitativos para la Investigación (Cátedra).}

\cventry{2021 -- 2022}{Universidad de Talca}
\cventrysub{Métodos Cualitativos para las Ciencias Sociales (Taller); Análisis de Datos para las Ciencias Sociales (Cátedra).}

\cventry{2022}{Universidad Adolfo Ibáñez}
\cventrysub{Redes Sociales, Democracia y Ciudadanía Digital (Taller).}

% ============================================
% 6. PONENCIAS (SELECCIÓN)
% ============================================
\section{Ponencias (Selección)}

\begin{itemize}[leftmargin=*]
    \item (2024) I Congreso Internacional de Evaluación de la Gobernanza. Universidad Católica de Temuco, Chile.
    \item (2024) I Congreso RELASSC. Recife, Brasil.
    \item (2023) XIV Congreso Chileno de Ciencia Política. Universidad de Santiago, Chile.
    \item (2023) VIII Congreso Uruguayo de Ciencia Política. Universidad Católica del Uruguay.
    \item (2022) IV Congreso Internacional ICCA. Universidad Complutense de Madrid, España.
    \item (2022) V Congreso Latinoamericano FLACSO. Montevideo, Uruguay.
\end{itemize}

% ============================================
% 7. BECAS Y DISTINCIONES
% ============================================
\section{Becas, Distinciones y Fondos}

\begin{itemize}[leftmargin=*]
    \item \textbf{2023--2024}: Beca Pasantía Internacional ANID (12,000 USD). IIGG, Universidad de Buenos Aires.
    \item \textbf{2021--2022}: Beca Beneficios Complementarios ANID (6,000 USD).
    \item \textbf{2020}: Beca Doctorado Nacional ANID (8,000 USD).
    \item \textbf{2019}: Fondo de Desarrollo Institucional, Ministerio de Educación (7,000 USD).
    \item \textbf{2018}: Mejor Egresado de la Generación, Universidad de Talca.
\end{itemize}

% ============================================
% 8. HABILIDADES
% ============================================
\section{Software e Idiomas}

\textbf{Programación:} R (Avanzado), Python (Intermedio), SQL (Intermedio), \LaTeX\ (Avanzado).\\
\textbf{Análisis cuantitativo:} SPSS, STATA, R, Jamovi (Avanzado).\\
\textbf{Análisis geoespacial:} QGIS (Avanzado).\\
\textbf{Análisis cualitativo:} ATLAS.ti (Avanzado), VOSviewer (Intermedio).\\
\textbf{Idiomas:} Español (Nativo), Inglés (Intermedio escrito, Básico oral).

% ============================================
% 9. ÁRBITRO EN REVISTAS
% ============================================
\section{Árbitro en Revistas Académicas}

\textit{Ethnic and Racial Studies}; \textit{Journal of Infrastructure, Policy and Development}; \textit{PLOSone}; \textit{Economía y Política}; \textit{Economies}; \textit{Sustainability}; \textit{Revista Desarrollo y Sociedad}; \textit{Revista Iberoamericana de Estudios Municipales (RIEM)}.

% ============================================
% 10. REFERENCIAS
% ============================================
\section{Referencias}

\textbf{Mauricio Morales}. Universidad de Talca. Doctor en Ciencia Política, PUC. \href{mailto:mmoralesq@utalca.cl}{mmoralesq@utalca.cl}

\textbf{Aldo Mascareño}. Universidad Adolfo Ibáñez. Doctor en Sociología, U. Bielefeld. \href{mailto:amascareno@cepchile.cl}{amascareno@cepchile.cl}

\textbf{Bernardita Escobar}. Universidad de Valparaíso. Doctora en Economía, Cambridge. \href{mailto:beni.escobar.andrae@gmail.com}{beni.escobar.andrae@gmail.com}

\textbf{Daniel Chernilo}. Universidad Adolfo Ibáñez. Doctor en Sociología, U. Warwick. \href{mailto:daniel.chernilo.s@uai.cl}{daniel.chernilo.s@uai.cl}

\end{document}
